\section{Connectivity}

\entry{SCC (Tarjan, iterative)}{$O(n + m)$ \quad mem $n + m$}{%
Directed. Takes any adjacency of an integer type. No recursion, so depth never matters;
self-loops and multi-edges are fine.\\[0.3ex]
\code{comp} is in \textbf{topological} order --- every edge \code{u -> v} has
\code{comp[u] <= comp[v]} --- so a DP over the condensation is just a loop over
\code{0 .. nc-1}, no toposort. \code{members()} gives the vertices of each comp,
\code{dag()} the condensation with duplicates removed.\\[0.3ex]
A comp holds a cycle iff it has $\ge 2$ vertices, or 1 with a self-loop; the graph is a
DAG iff \code{nc == n} and no self-loop exists. Sources of the condensation have
in-degree 0 in \code{dag()}, sinks out-degree 0 --- count degrees, \code{comp[v] == 0}
is not enough.\\[0.3ex]
\textbf{2-SAT.} Node \code{2i} is $x_i$ false, \code{2i+1} is true, so \code{n = 2*vars}.
A clause $(a \lor b)$ adds \emph{both} \code{!a -> b} and \code{!b -> a}; forcing $a$
true is the single edge \code{!a -> a}. Satisfiable iff
\code{comp[2i] != comp[2i+1]} for every $i$, and then
$x_i$ = \code{comp[2i+1] > comp[2i]} (needs \code{comp} topological, as here).}

%LABEL SCC
\begin{lstlisting}
struct SCC {
    int n, nc = 0;              // nc = number of comps
    vector<vector<int>> g;
    vector<int> comp;           // topological order
    template<class G>          // any vector<vector<intlike>>
    SCC(const G& adj) : n(adj.size()), g(n), comp(n, -1) {
        for (int u = 0; u < n; u++)
            for (auto v : adj[u]) g[u].push_back((int)v);
        vector<int> tin(n, -1), low(n), ptr(n, 0), stk, cs;
        vector<char> on(n, 0);
        int timer = 0;
        for (int r = 0; r < n; r++) {
            if (tin[r] != -1) continue;
            tin[r] = low[r] = timer++;
            stk.push_back(r); on[r] = 1; cs.push_back(r);
            while (!cs.empty()) {
                int u = cs.back();
                if (ptr[u] < (int)g[u].size()) {
                    int v = g[u][ptr[u]++];
                    if (tin[v] == -1) {
                        tin[v] = low[v] = timer++;
                        stk.push_back(v); on[v] = 1;
                        cs.push_back(v);
                    } else if (on[v])
                        low[u] = min(low[u], tin[v]);
                } else {
                    cs.pop_back();
                    if (!cs.empty()) {
                        int p = cs.back();
                        low[p] = min(low[p], low[u]);
                    }
                    if (low[u] == tin[u]) {  // u roots comp
                        int w;
                        do {
                            w = stk.back(); stk.pop_back();
                            on[w] = 0; comp[w] = nc;
                        } while (w != u);
                        nc++;
                    }
                }
            }
        }
        for (int v = 0; v < n; v++)         // -> topological
            comp[v] = nc - 1 - comp[v];
    }
    vector<vector<int>> members() {
        vector<vector<int>> c(nc);
        for (int v = 0; v < n; v++) c[comp[v]].push_back(v);
        return c;
    }
    vector<vector<int>> dag() {      // condensation, dedup
        vector<vector<int>> d(nc);
        for (int u = 0; u < n; u++)
            for (int v : g[u])
                if (comp[u] != comp[v])
                    d[comp[u]].push_back(comp[v]);
        for (auto& a : d) {
            sort(all(a));
            a.erase(unique(all(a)), a.end());
        }
        return d;
    }
};
\end{lstlisting}

\entry{Bridges, cut vertices, 2ECC}{$O(n + m)$ \quad mem $n + m$}{%
Undirected, given as an \textbf{edge list} (pairs, 2-element arrays or tuples of any
integer type), so \code{bri[i]} indexes your own input. Multi-edges and self-loops are
safe, disconnected input is fine, no recursion.\\[0.3ex]
\begin{ctab}{0.15}{0.80}
\code{bri[i]} & input edge \code{i} is a bridge.\\
\code{art[v]} & \code{v} is a cut vertex.\\
\code{comp[v]} & 2-edge-connected comp: all you reach without crossing a bridge.\\
\code{tree()} & the bridge forest on those comps.\\
\end{ctab}\\[0.3ex]
Bridges are exactly the edges between comps, so \code{tree()} is a forest, one tree per
connected component --- feed it to \code{LCA} / \code{HLD} for ``how many bridges on
the \code{u}--\code{v} path''. Adding one edge kills every bridge on the path joining
its two comps.\\[0.3ex]
The parent edge is skipped by \textbf{edge id}, never by vertex --- by vertex reports
parallel edges as bridges, which is the usual wrong answer. A bridge endpoint need not
be a cut vertex (degree 1) and a cut vertex need not carry one (two cycles sharing it);
the DFS root is one iff it has $\ge 2$ children.}

\begin{lstlisting}
struct Bridges {
    int n, m, nc = 0;                  // nc = number of 2ecc
    vector<vector<pair<int,int>>> g;   // (to, edge id)
    vector<int> comp;                  // 2ecc id per vertex
    vector<char> bri, art;             // per edge / vertex
    vector<pair<int,int>> es;
    template<class E>      // list of pair / array / tuple
    Bridges(int n, const E& edges)
            : n(n), m(edges.size()), g(n), comp(n, -1),
              bri(m, 0), art(n, 0) {
        es.reserve(m);
        for (auto& e : edges)
            es.push_back({(int)get<0>(e), (int)get<1>(e)});
        vector<int> deg(n, 0);
        for (int i = 0; i < m; i++) {
            deg[es[i].first]++; deg[es[i].second]++;
        }
        for (int v = 0; v < n; v++) g[v].reserve(deg[v]);
        for (int i = 0; i < m; i++) {
            g[es[i].first].push_back({es[i].second, i});
            g[es[i].second].push_back({es[i].first, i});
        }
        vector<int> tin(n, -1), low(n), ptr(n, 0), cs;
        vector<int> pe(n, -1);
        int timer = 0;
        for (int r = 0; r < n; r++) {
            if (tin[r] != -1) continue;
            int kids = 0;
            tin[r] = low[r] = timer++; cs.push_back(r);
            while (!cs.empty()) {
                int u = cs.back();
                if (ptr[u] < (int)g[u].size()) {
                    auto [v, id] = g[u][ptr[u]++];
                    if (id == pe[u]) continue;   // skip by ID
                    if (tin[v] == -1) {
                        pe[v] = id; tin[v] = low[v] = timer++;
                        cs.push_back(v); kids += (u == r);
                    } else low[u] = min(low[u], tin[v]);
                } else {
                    cs.pop_back();
                    if (cs.empty()) break;
                    int p = cs.back();
                    low[p] = min(low[p], low[u]);
                    if (low[u] > tin[p]) bri[pe[u]] = 1;
                    if (low[u] >= tin[p]) art[p] = 1;
                }
            }
            art[r] = kids >= 2;                  // root rule
        }
        vector<int> q;                // 2ecc: drop bridges
        for (int r = 0; r < n; r++) {
            if (comp[r] != -1) continue;
            comp[r] = nc; q.clear(); q.push_back(r);
            while (!q.empty()) {
                int u = q.back(); q.pop_back();
                for (auto [v, id] : g[u])
                    if (!bri[id] && comp[v] == -1) {
                        comp[v] = nc; q.push_back(v);
                    }
            }
            nc++;
        }
    }
    Bridges(int n, initializer_list<pair<int,int>> e)
            : Bridges(n, vector<pair<int,int>>(e)) {}
    vector<vector<int>> tree() {     // bridge forest
        vector<vector<int>> t(nc);
        for (int i = 0; i < m; i++) if (bri[i]) {
            int a = comp[es[i].first], b = comp[es[i].second];
            t[a].push_back(b); t[b].push_back(a);
        }
        return t;
    }
};
\end{lstlisting}

\entry{Topological sort (Kahn)}{$O(n + m)$ \quad mem $n$}{%
Directed. Returns an empty vector iff the graph has a cycle, so it doubles as a cycle
test. If you already built \code{SCC} you have a topological order for free in
\code{comp} --- use this only when you do not need the components.\\[0.3ex]
Pop from the back instead of the front and you get a DFS-flavoured order; use a
\code{priority\_queue} for the lexicographically smallest one.}

\begin{lstlisting}
vector<int> topo(const vector<vector<int>>& g) {
    int n = (int)g.size();
    vector<int> deg(n, 0), q;
    for (int u = 0; u < n; u++)
        for (int v : g[u]) deg[v]++;
    for (int u = 0; u < n; u++) if (!deg[u]) q.push_back(u);
    for (int i = 0; i < (int)q.size(); i++)
        for (int v : g[q[i]])
            if (--deg[v] == 0) q.push_back(v);
    return (int)q.size() == n ? q : vector<int>{};
}
\end{lstlisting}

\entry{2-SAT}{KACTL \quad $O(N + E)$}{%
The ready-made version of the 2-SAT note under \emph{SCC}. Recognise it: every item
has \textbf{two choices} (on/off, left/right, colour A/B) and every constraint mentions
\textbf{at most two} items (``not both'', ``at least one of'', ``same'', ``if this then
that'').\\[0.3ex]
Variables are \code{0..N-1}; \code{\textasciitilde x} is ``x is false''.
\code{either(a, b)} is $a \lor b$, so: implication $a \Rightarrow b$ is
\code{either(\textasciitilde a, b)}; not both is
\code{either(\textasciitilde a, \textasciitilde b)}; $a = b$ is
\code{either(\textasciitilde a, b)} $+$ \code{either(a, \textasciitilde b)}.
\code{setValue(x)} forces x true. \code{atMostOne(\{..\})} adds helper variables so it
stays linear --- do not write the $O(k^2)$ pairs.\\[0.3ex]
\code{solve()} returns satisfiable; \code{values[i]} is then an assignment. Recursive
DFS: depth up to $2N$, fine to $\sim10^5$ variables.}

%LABEL twosat
\begin{lstlisting}
struct TwoSat {
    int N;
    vector<vi> gr;
    vi values; // 0 = false, 1 = true
    TwoSat(int n = 0) : N(n), gr(2*n) {}
    int addVar() { // (optional)
        gr.emplace_back();
        gr.emplace_back();
        return N++;
    }
    void either(int f, int j) {
        f = max(2*f, -1-2*f);
        j = max(2*j, -1-2*j);
        gr[f].push_back(j^1);
        gr[j].push_back(f^1);
    }
    void setValue(int x) { either(x, x); }
    void atMostOne(const vi& li) { // (optional)
        if (li.size() <= 1) return;
        int cur = ~li[0];
        rep(i,2,(ll)li.size()) {
            int next = addVar();
            either(cur, ~li[i]);
            either(cur, next);
            either(~li[i], next);
            cur = ~next;
        }
        either(cur, ~li[1]);
    }
    vector<int> val, comp, z; int time = 0;
    int dfs(int i) {
        int low = val[i] = ++time, x; z.push_back(i);
        for(int e : gr[i]) if (!comp[e])
            low = min(low, val[e] ?: dfs(e));
        if (low == val[i]) do {
            x = z.back(); z.pop_back();
            comp[x] = low;
            if (values[x>>1] == -1)
                values[x>>1] = x&1;
        } while (x != i);
        return val[i] = low;
    }
    bool solve() {
        values.assign(N, -1);
        val.assign(2*N, 0); comp = val;
        rep(i,0,2*N) if (!comp[i]) dfs(i);
        rep(i,0,N) if (comp[2*i] == comp[2*i+1]) return 0;
        return 1;
    }
};
\end{lstlisting}

\entry{Maximum clique}{KACTL \quad $\sim$1 s for $n = 155$, dense}{%
Largest set of pairwise adjacent vertices, exponential but heavily pruned (greedy
colouring bounds). Input: symmetric \code{bitset} adjacency, \textbf{no self-loops};
raise \code{200} to your $n$. Returns the vertices.\\[0.3ex]
\textbf{Maximum independent set} of a general graph = maximum clique of the
\emph{complement} (flip every off-diagonal bit). Bipartite graphs do not need this ---
see the vertex cover entry under matching. Typical bounds: $n \le 40$--$150$, ``pick as
many as possible, no two conflicting''.\\[0.3ex]
KACTL uses \code{sz} everywhere; it is defined and undefined around the struct so
\code{Treap} still compiles.}

%LABEL clique
\begin{lstlisting}
#define sz(x) (int)(x).size()
typedef vector<bitset<200>> vb;
struct Maxclique {
    double limit=0.025, pk=0;
    struct Vertex { int i, d=0; };
    typedef vector<Vertex> vv;
    vb e;
    vv V;
    vector<vi> C;
    vi qmax, q, S, old;
    void init(vv& r) {
        for (auto& v : r) v.d = 0;
        for (auto& v : r) for (auto j : r) v.d += e[v.i][j.i];
        sort(all(r), [](auto a, auto b) { return a.d>b.d; });
        int mxD = r[0].d;
        rep(i,0,sz(r)) r[i].d = min((int)i, mxD) + 1;
    }
    void expand(vv& R, int lev = 1) {
        S[lev] += S[lev - 1] - old[lev];
        old[lev] = S[lev - 1];
        while (sz(R)) {
            if (sz(q) + R.back().d <= sz(qmax)) return;
            q.push_back(R.back().i);
            vv T;
            for(auto v:R) if (e[R.back().i][v.i])
                T.push_back({v.i});
            if (sz(T)) {
                if (S[lev]++ / ++pk < limit) init(T);
                int j = 0, mxk = 1;
                int mnk = max(sz(qmax) - sz(q) + 1, 1);
                C[1].clear(), C[2].clear();
                for (auto v : T) {
                    int k = 1;
                    auto f = [&](int i) { return e[v.i][i]; };
                    while (any_of(all(C[k]), f)) k++;
                    if (k > mxk) mxk = k, C[mxk + 1].clear();
                    if (k < mnk) T[j++].i = v.i;
                    C[k].push_back(v.i);
                }
                if (j > 0) T[j - 1].d = 0;
                rep(k,mnk,mxk + 1) for (int i : C[k])
                    T[j].i = i, T[j++].d = k;
                expand(T, lev + 1);
            } else if (sz(q) > sz(qmax)) qmax = q;
            q.pop_back(), R.pop_back();
        }
    }
    vi maxClique() { init(V), expand(V); return qmax; }
    Maxclique(vb conn) : e(conn), C(sz(e)+1), S(sz(C)),
        old(S) {
        rep(i,0,sz(e)) V.push_back({(int)i});
    }
};
#undef sz
\end{lstlisting}
